\section{无穷小量与无穷大量的比较}\label{sec:3.5}

先看数列的情形．设 $x_n$，$y_n$ 都是无穷小量，即 $\lim_{n \to \infty} x_n = 0$，$\lim_{n \to \infty} y_n = 0$．这时数列 $\frac{x_n}{y_n}$ 的趋向，对不同的 $x_n$，$y_n$，可能出现各种情形．

\begin{example}\label{ex:3-5-1}
$x_n = \frac{cn}{n^2+1}$，$y_n = \frac{1}{n}$，其中 $c \neq 0$ 是常数．这时
\[
\lim_{n \to \infty} \frac{x_n}{y_n} = \lim_{n \to \infty} \frac{cn^2}{n^2+1} = c,
\]
即 $\frac{x_n}{y_n}$ 有非 $0$ 的极限．
\end{example}

\begin{example}\label{ex:3-5-2}
$x_n = \frac{1}{n^2}$，$y_n = \frac{1}{n}$，则
\[
\lim_{n \to \infty} \frac{x_n}{y_n} = \lim_{n \to \infty} \frac{n}{n^2} = 0,
\]
或等价于 $\lim_{n \to \infty} \frac{y_n}{x_n} = \infty$．
\end{example}

\begin{example}\label{ex:3-5-3}
$x_n = \frac{(-1)^n}{n}$，$y_n = \frac{1}{n}$，这时数列 $\frac{x_n}{y_n} = (-1)^n$ 没有极限，但其绝对值 $\left|\frac{x_n}{y_n}\right| = 1$ 有正的下界与上界．
\end{example}

\begin{example}\label{ex:3-5-4}
$x_n = \frac{1+(-1)^n}{n}$，$y_n = \frac{1}{n^2}$，这时 $\frac{x_n}{y_n} = [1+(-1)^n]n$，写出来就是：$0, 4, 0, 8, 0, 12, 0, 16, \cdots$，它既没有极限，又没有上界与正的下界．
\end{example}

比较这几种情形，可得下面的定义：

\begin{definition}\label{def:3-11}
设 $\lim_{n \to \infty} x_n = 0$，$\lim_{n \to \infty} y_n = 0$．
\begin{enumerate}[label=(\arabic*)]
\item 若存在 $A > 0$，$B > 0$ 及正整数 $N$，使得当 $n > N$ 时，有
\[
0 < A \leqslant \frac{|x_n|}{|y_n|} \leqslant B,
\]
则称 $x_n$ 与 $y_n$ 是同阶的无穷小量．
\item 若
\[
\lim_{n \to \infty} \frac{x_n}{y_n} = 1,
\]
则称 $x_n$ 与 $y_n$ 为等价的无穷小量，记为 $x_n \sim y_n$ ($n \to \infty$)．
\item 若
\[
\lim_{n \to \infty} \frac{x_n}{y_n} = 0,
\]
则称 $x_n$ 为较 $y_n$ 高阶的无穷小量，或称 $y_n$ 是较 $x_n$ 低阶的无穷小量，记为
\[
x_n = o(y_n) \quad (n \to \infty).
\]
\end{enumerate}
\end{definition}

显然，$x_n = o(y_n)$ 是说，$x_n = \alpha_n y_n$，而 $\lim_{n \to \infty} \alpha_n = 0$，因此对任意给定 $\varepsilon > 0$，存在 $N$，当 $n > N$ 时，有
\[
|x_n| < \varepsilon |y_n|,
\]
这说明 $x_n$ 比 $y_n$ 趋向于零要快得多．

自然地，符号 $x_n = o(1)$，就表示 $x_n$ 是无穷小量．

当 $x_n$ 与 $y_n$ 是同阶无穷小量时，则当 $n$ 充分大后，其绝对值互相被一常数倍限制着，即
\[
|x_n| \leqslant B|y_n|, \quad |y_n| \leqslant \frac{1}{A}|x_n|,
\]
也就是说，它们趋向于 $0$ 的速度可以用常数倍来度量，不像高阶无穷小量那样，速度快一个``量级''．如果
\[
\lim_{n \to \infty} \frac{x_n}{y_n} = a \neq 0,
\]
则 $x_n$ 与 $y_n$ 是同阶的，这是因为由极限的性质知
\[
\lim_{n \to \infty} \frac{|x_n|}{|y_n|} = |a| > 0,
\]
因此，存在 $N$，当 $n > N$ 时，有
\[
0 < \frac{|a|}{2} \leqslant \frac{|x_n|}{|y_n|} \leqslant \frac{3}{2}|a|.
\]

当 $x_n$ 与 $y_n$ 是等价的无穷小量时，则对任意给定的 $\varepsilon > 0$，存在 $N$，当 $n > N$ 时，有
\[
1 - \varepsilon < \frac{x_n}{y_n} < 1 + \varepsilon,
\]
这说明 $x_n$ 与 $y_n$ 不仅同阶，而且当 $n$ 很大时，它们差不多是相等的．进一步，由 $\lim_{n \to \infty} \frac{x_n}{y_n} = 1$，记 $\frac{x_n}{y_n} - 1 = \alpha_n$，知 $\lim_{n \to \infty} \alpha_n = 0$．因此
\[
x_n - y_n = \alpha_n y_n = o(y_n),
\]
这表明 $x_n$ 与 $y_n$ 的差是一个较 $y_n$（也较 $x_n$）高阶的无穷小量．

回到本节开始时的例子，知 $\frac{cn}{n^2+1}$（$c \neq 0$）与 $\frac{1}{n}$ 是同阶的，$\frac{(-1)^n}{n}$ 与 $\frac{1}{n}$ 也是同阶的，$\frac{1}{n^2}$ 是较 $\frac{1}{n}$ 高阶的无穷小，而例~\ref{ex:3-5-4} 中的 $\frac{1+(-1)^n}{n}$ 与 $\frac{1}{n^2}$ 是无法比较的两个无穷小量．

我们还要引进一个记号：如果 $\frac{x_n}{y_n}$ 是有界的，即存在常数 $M > 0$ 使得
\[
\left|\frac{x_n}{y_n}\right| \leqslant M,
\]
那么记 $x_n = O(y_n)$．显然 $x_n = o(y_n)$ 蕴含了 $x_n = O(y_n)$．

如果选定 $\frac{1}{n}$ 作为无穷小量的标准，若 $x_n$ 满足
\[
\lim_{n \to \infty} \frac{x_n}{\frac{1}{n^\alpha}} = a \neq 0,
\]
其中 $\alpha$ 是某个正常数，则称 $x_n$ 是 $\alpha$ 阶的无穷小量，这时
\[
x_n = \frac{a}{n^\alpha} + \alpha_n,
\]
$\alpha_n$ 是较 $\frac{1}{n^\alpha}$ 高阶的无穷小量，我们把 $\frac{a}{n^\alpha}$ 称为 $x_n$ 的主部．例如 $\sqrt{n+1} - \sqrt{n} = \frac{1}{\sqrt{n+1} + \sqrt{n}}$ 是 $\frac{1}{2}$ 阶的，由于
\[
\frac{1}{\sqrt{n+1} + \sqrt{n}} - \frac{1}{2\sqrt{n}} = \frac{\sqrt{n} - \sqrt{n+1}}{2\sqrt{n}(\sqrt{n+1} + \sqrt{n})} = o\left(\frac{1}{\sqrt{n}}\right),
\]
故其主部是 $\frac{1}{2\sqrt{n}}$．

类似的概念也可转移到连续变量的函数极限．以 $x \to a$ 为例．设 $\lim_{x \to a} f(x) = 0$，$\lim_{x \to a} g(x) = 0$．若存在 $A > 0$，$B > 0$，以及 $\delta > 0$，当 $0 < |x - a| < \delta$ 时，有
\[
A \leqslant \left|\frac{f(x)}{g(x)}\right| \leqslant B,
\]
则称 $f(x)$ 与 $g(x)$ 在 $x \to a$ 时是同阶的无穷小量．若
\[
\lim_{x \to a} \frac{f(x)}{g(x)} = 1,
\]
则称 $f(x)$ 与 $g(x)$ 当 $x \to a$ 时是等价的无穷小量，记为
\[
f(x) \sim g(x) \quad (x \to a).
\]
若
\[
\lim_{x \to a} \frac{f(x)}{g(x)} = 0,
\]
则称当 $x \to a$ 时 $f(x)$ 是较 $g(x)$ 高阶的无穷小量，记为
\[
f(x) = o(g(x)) \quad (x \to a).
\]
若存在 $\delta > 0$，使得 $\frac{f(x)}{g(x)}$ 在 $0 < |x - a| < \delta$ 有界，则记为
\[
f(x) = O(g(x)) \quad (x \to a).
\]
类似于数列的情形，由 $f(x) = o(g(x)) \ (x \to a)$，可推得 $f(x) = O(g(x)) \ (x \to a)$．

根据以前计算过的极限，我们知道当 $x \to 0$ 时 $\sin x$ 等价于 $x$，$\ln(1+x)$ 等价于 $x$，等等．

利用等价无穷小量，可以简化求极限的运算．

\begin{example}\label{ex:3-5-5}
求极限 $\lim_{x \to 0} \frac{\ln(1+2x^2)}{\tan^2 x}$．
\end{example}

\begin{proof}[解]
因为
\[
\ln(1+2x^2) \sim 2x^2 \quad (x \to 0),
\]
\[
\tan x \sim x \quad (x \to 0),
\]
所以
\[
\lim_{x \to 0} \frac{\ln(1+2x^2)}{\tan^2 x} = \lim_{x \to 0} \frac{2x^2}{x^2} = 2.
\]
\end{proof}

总起来说，无穷小量的比较是用它们趋向于 $0$ 的速度快慢来衡量的．这个观念是分析中十分重要的观念，现在还不能很好体会．以后学多了，便会逐步加深理解．

无穷小量的阶是无止境的．换句话说，任一无穷小量都存在比它更高阶的无穷小量，也存在比它更低阶的无穷小量．例如，下面的数列中，后者是比前者更高阶的无穷小量：
\[
\cdots, \left\{\frac{1}{\ln\ln n}\right\}, \left\{\frac{1}{\ln n}\right\}, \left\{\frac{1}{n}\right\}, \left\{\frac{1}{e^n}\right\}, \left\{\frac{1}{n!}\right\}, \left\{\frac{1}{n^n}\right\}, \cdots
\]
不仅如此，在上面相邻的两个无穷小量之间还可插入无穷小量，使得后者是比前者更高阶的无穷小量．如 $\left\{\frac{1}{n}\right\}$ 和 $\left\{\frac{1}{e^n}\right\}$ 之间可插入无穷小量 $\left\{\frac{1}{n^2}\right\}$；$\left\{\frac{1}{e^n}\right\}$ 和 $\left\{\frac{1}{n!}\right\}$ 之间可插入无穷小量 $\left\{\frac{1}{3^n}\right\}$ 等等．

根据无穷大量与无穷小量互为倒数的关系，当 $x \to x_0$ 时，若 $f(x)$ 是比 $g(x)$ 更高阶的无穷小量，则 $\frac{1}{f(x)}$ 是比 $\frac{1}{g(x)}$ 更高阶的无穷大量．因此完全类似地可以讨论无穷大量的比较，在这里不再叙述了．

\begin{xiti}
\item 当 $x \to 0$ 时，以 $x$ 为标准求下列无穷小量的阶：
\begin{ziti}
\item $\sin 2x - 2\sin x$；
\item $\frac{1}{1+x} - (1 - x)$；
\item $\sqrt{1 + \tan x} - \sqrt{1 - \sin x}$；
\item $\ln(1+x)$；
\item $\sqrt[n]{1+x} - 1$；
\item $e^x - 1$．
\end{ziti}

\item 当 $x \to \pm \infty$ 时，以 $x$ 为标准求下列无穷大量的阶：
\begin{ziti}
\item $x^2 + x^6$；
\item $4x^2 + 6x^4 - x^5$；
\item $\sqrt[3]{x^2\sin\frac{1}{x}}$；
\item $\frac{x^3+1}{x^2+2x-3}$；
\item $x^2 \arctan\frac{1}{x}$．
\end{ziti}

\item 当 $x \to 0$ 时，下列等式成立吗？
\begin{ziti}
\item $o(x^2) = o(x)$；
\item $O(x^2) = o(x)$；
\item $x \cdot o(x^2) = o(x^3)$；
\item $\frac{o(x^2)}{x} = o(x)$；
\item $\frac{o(x^2)}{o(x)} = o(x)$；
\item $o(x) = O(x^2)$．
\end{ziti}

\item 试证下列各题：
\begin{ziti}
\item $x\sin\sqrt{x} = O(x^{3/2})$ ($x \to 0^+$)；
\item $2x^3 + 2x^2 = O(x^3)$ ($x \to \infty$)；
\item $o(g(x)) \pm o(g(x)) = o(g(x))$ ($x \to x_0$)；
\item $o(x^m) + o(x^n) = o(x^n)$ ($x \to 0$)，$m > n > 0$；
\item $o(x^m)o(x^n) = o(x^{m+n})$ ($x \to 0$)，$m, n > 0$．
\end{ziti}

\item 证明下列各式：
\begin{ziti}
\item $\tan x \sim x$ ($x \to 0$)；
\item $\arcsin x \sim x$ ($x \to 0$)；
\item $\arctan x \sim x$ ($x \to 0$)；
\item $1 - \cos x \sim \frac{1}{2}x^2$ ($x \to 0$)；
\item $e^x - 1 \sim x$ ($x \to 0$)；
\item $(1 + x)^\alpha - 1 \sim \alpha x$ ($x \to 0$)，其中 $\alpha \neq 0$．
\end{ziti}

\item 运用等价无穷小量求极限：
\begin{ziti}
\item $\lim_{x \to \infty} \frac{2\arctan\frac{1}{x}}{x - \cos x}$；
\item $\lim_{x \to 0} \frac{\sqrt{1+x^2} - 1}{1 - \cos x}$；
\item $\lim_{x \to 0} \frac{x\ln(1+x)}{\sin x^2}$；
\item $\lim_{x \to 0} \frac{e^{x^2} - 1}{x\sin x}$．
\end{ziti}

\item 设 $f(x) \sim g(x)$ ($x \to x_0$)，证明：
\[
f(x) - g(x) = o(f(x)) \quad \text{与} \quad f(x) - g(x) = o(g(x)).
\]

\item 设 $x \to a$ 时，$f_1(x)$ 与 $f_2(x)$ 为等价无穷小，$g_1(x)$ 与 $g_2(x)$ 是等价无穷大，且 $\lim_{x \to a} f_2(x)g_2(x)$ 存在，求证：
\[
\lim_{x \to a} f_1(x)g_1(x) = \lim_{x \to a} f_2(x)g_2(x).
\]
\end{xiti}
